We consider a single forecast of length H for the following complexity analysis, with a \NBEATS\ and a \ours\ architecture of $B$ blocks. We do not consider the batch dimension. We consider most practical situations, the input size $L=\mathcal{O}(H)$ linked to the horizon length.

The block operation described by Equation~(\ref{equation:nhits_projections}) has complexity dominated by the fully connected layers of $\mathcal{O}(H\,N_{h})$, with $N_{h}$ the number of hidden units that we treat as a constant. The depth of stacked blocks in the \NBEATSg\ architecture, that endows it with its expressivity, is associated to a computational complexity that scales linearly $\mathcal{O}(H B)$, with $B$ the number of blocks.

The block operation described by Equation~(\ref{equation:nhits_projections}) has complexity dominated by the fully connected layers of $\mathcal{O}(H\,N_{h})$, with $N_{h}$ the number of hidden units that we treat as a constant. The depth of stacked blocks in the \NBEATSg\ architecture, which endows it with its expressivity, is associated with a computational complexity that scales linearly $\mathcal{O}(H B)$, with $B$ the number of blocks.

In contrast the \ours\ architecture that specializes each stack in different frequencies, through the expressivity ratios, can greatly reduce the amount of parameters needed for each layer. When we use \emph{exponentially increasing expressivity} ratios through the depth of the architecture blocks it allows to model complex dependencies, while controlling the number of parameters used on each output layer. If the \emph{expressivity ratio} is defined as $r_{\ell} = r^{l}$ then the space complexity of \ours\ scales geometrically $\mathcal{O}(\sum^{B}_{l=0} H r^{l}) =\mathcal{O}\left((H(1-r^{B})/(1-r)\right)$.

\begin{table}[!ht] %!ht ht!
%\scriptsize
\footnotesize
\centering
\tiny
\caption[caption]{Computational complexity of neural based forecasting methods as a function of the output size $H$. For simplicity, we assume that the input size $L$ scales linearly with respect to $H$. For \ours\ and \NBEATS\ we also consider the network's $B$ blocks.} %\\\hspace{\textwidth} \tiny{\textsuperscript{*} Reported complexity analysis by benchmark methods.}
\begin{tabular}{l | cc} \toprule
\textsc{Model}        & \textsc{Time}                   & \textsc{Memory}       \\ \midrule
\LSTM                 & $\mathcal{O}(H)$                & $O(H)$                \\
\ESRNN                & $\mathcal{O}(H)$                & $O(H)$                \\
\TCN                  & $\mathcal{O}(H)$                & $O(H)$                \\ 
\Transformer          & $\mathcal{O}(H^2)$              & $O(H^2)$              \\
\Reformer             & $\mathcal{O}(H\,log H)$         & $O(H\,log H)$         \\
\Informer             & $\mathcal{O}(H\,log H)$         & $O(H\,log H)$         \\
\Autoformer           & $\mathcal{O}(H\,log H)$         & $O(H\,log H)$         \\ 
\LogTrans             & $\mathcal{O}(H\,log H)$         & $O(H^2)$              \\ \midrule
\NBEATSi              & $\mathcal{O}(H^2B)$             & $O(H^2B)$             \\
\NBEATSg              & $\mathcal{O}(HB)$               & $O(HB)$               \\ 
\ours                 & $\mathcal{O}(H(1-r^{B})/(1-r))$ & $O(H(1-r^{B})/(1-r))$ \\ \bottomrule
\end{tabular}
\end{table}